
\documentclass{article}

\usepackage[left=3cm,top=3cm,right=2cm,bottom=2cm,a4paper]{geometry}
\usepackage{anyfontsize}
\usepackage[indent=0cm,skip=.1cm]{parskip}
\usepackage[brazil]{babel}

\begin{document}
\fontsize{12pt}{18pt}\selectfont

%=====================================
% CAPA
%=====================================

\begin{titlepage}

\begin{center}

FACULDADE MAURÍCIO DE NASSAU

\vspace{2cm}

Phellipe de Amorim martins - 01782567

Anthony Gabriel Lemos Baracho - 01813891

Ruan Kelvin Vieira dos Santos - 01797916

Everton Hasabias Furtunato Celestino - 1796727

Adriano Silva do Nascimento - 01769583

Fernanda Nogueira de França - 01735985

\vfill

{\fontsize{16pt}{20pt}\selectfont\textbf{CÓDIGOS DE ALTA PERFORMANCE}}

\vspace{0.5cm}

NEXPAY

\vfill

Maceió - AL

Junho de 2026

\end{center}

\end{titlepage}

%=====================================
% FOLHA DE ROSTO
%=====================================

\begin{titlepage}

\begin{center}

Phellipe de Amorim martins - 01782567

Anthony Gabriel Lemos Baracho - 01813891

Ruan Kelvin Vieira dos Santos - 01797916

Everton Hasabias Furtunato Celestino - 1796727

Adriano Silva do Nascimento - 01769583

Fernanda Nogueira de França - 01735985

\vspace{6cm}

{\fontsize{16pt}{20pt}\selectfont\textbf{CÓDIGOS DE ALTA PERFORMANCE}}

\vspace{0.5cm}

NEXPAY

\end{center}

\vspace{3cm}

\hspace*{7cm}
\parbox{8cm}{
Trabalho apresentado à disciplina de
Códigos de Alta perfomance do curso de
Ciências da Computação da Faculdade Maurício de Nassau,
como requisito para obtenção de nota.

\vspace{0.5cm}

Professor(a): Anderlan Henrique Batista Siqueira
}

\vfill

\begin{center}

Maceió - AL

Junho de 2026

\end{center}

\end{titlepage}


\tableofcontents
\newpag

\clearpage
\section{Introdução ao sistema}

\subsection{Nome do Sistema}

Nexpay.

\subsection{Problema}

Ajudar pessoas que possuem dificuldades para organizar e controlar suas finanças pessoais.

\subsection{Público-Alvo}

Jovens e adultos que desejam melhorar o gerenciamento de seus recursos financeiros.

\subsection{Objetivo Principal}

Centralizar a gestão financeira em uma única plataforma, permitindo ao usuário acompanhar receitas, despesas e saldo disponível de forma prática e segura.

\subsection{Contexto de Uso}

O sistema foi pensado principalmente para jovens que estão iniciando sua carreira profissional e ainda não possuem experiência suficiente para administrar seu dinheiro de maneira eficiente.

\section{Arquitetura como Decisão}

\subsection{Qual será a estrutura inicial do sistema?}

A estrutura inicial adotada será a arquitetura monolítica. Essa abordagem foi escolhida por ser mais simples de desenvolver, implantar e manter durante as fases iniciais do projeto, além de ser adequada para sistemas de pequeno porte.


\subsection{O que foi priorizado no início?}

Os aspectos priorizados na fase inicial do desenvolvimento foram:

\begin{itemize}
    \item Segurança dos dados;
    \item Estabilidade do sistema;
\end{itemize}

\subsection{O que foi deixado para uma evolução futura?}

As seguintes funcionalidades poderão ser implementadas em versões futuras:

\begin{itemize}
    \item Notificações inteligentes;
    \item Relatórios financeiros detalhados;
    \item Sistema de gestão avançado para Pessoa Jurídica (CNPJ).
\end{itemize}

\section{Requisitos}

\subsection{Requisitos Funcionais}

\begin{itemize}
    \item Registro de receitas financeiras;
    \item Registro de despesas financeiras;
    \item Edição e exclusão de movimentações financeiras;
    \item Categorização de despesas (alimentação, transporte, lazer, salário, entre outras);
    \item Exibição do saldo total atualizado do usuário;
\end{itemize}

\subsection{Requisitos Não Funcionais}

\begin{itemize}
    \item Segurança;
    \item Desempenho;
    \item Escalabilidade;
    \item Confiabilidade;
    \item Usabilidade.
\end{itemize}


\section{Tecnologias Utilizadas}

O desenvolvimento do sistema Nexpay utiliza tecnologias modernas que permitem a construção de uma aplicação web eficiente, escalável e de fácil manutenção.

\subsection{React}

O React é uma biblioteca JavaScript utilizada para a construção da interface do usuário. Sua arquitetura baseada em componentes facilita a reutilização de código, a organização do projeto e a manutenção da aplicação. Além disso, o React proporciona uma experiência mais dinâmica e responsiva para os usuários.

\subsection{Node.js}

O Node.js é o ambiente de execução utilizado no lado do servidor. Sua principal função é processar as requisições da aplicação, implementar as regras de negócio e realizar a comunicação com o banco de dados. A utilização do Node.js permite alto desempenho e facilita o desenvolvimento utilizando a linguagem JavaScript tanto no front-end quanto no back-end.

\subsection{Tailwind CSS}

O Tailwind CSS é um framework de estilização baseado em classes utilitárias. Ele foi escolhido por possibilitar a criação de interfaces modernas e responsivas de forma rápida e padronizada. Além disso, reduz a necessidade de escrever grandes arquivos CSS, tornando o desenvolvimento mais produtivo e organizado.



\section{Estrutura do Projeto}

O projeto Nexpay foi organizado seguindo uma estrutura modular, com separação das responsabilidades entre interface, gerenciamento de estado e páginas da aplicação. Essa organização facilita a manutenção, a escalabilidade e o desenvolvimento de novas funcionalidades.

\subsection{Estrutura de Diretórios}

\begin{itemize}
    \item \textbf{node\_modules/}: diretório gerado automaticamente pelo NPM, contendo todas as dependências utilizadas pelo projeto.

    \item \textbf{public/}: armazena arquivos públicos que podem ser acessados diretamente pela aplicação, como ícones, imagens e outros recursos estáticos.

    \item \textbf{src/}: diretório principal do código-fonte da aplicação.
    
    \begin{itemize}
        \item \textbf{assets/}: contém recursos estáticos utilizados pela interface, como imagens, ícones e arquivos de mídia.
        
        \item \textbf{context/}: responsável pelo gerenciamento de estado global da aplicação.
        
        \begin{itemize}
            \item \textbf{FinanceContext.jsx}: centraliza informações financeiras compartilhadas entre diferentes componentes da aplicação.
        \end{itemize}
        
        \item \textbf{pages/}: contém as páginas principais do sistema.
        
        \begin{itemize}
            \item \textbf{Home.jsx}: página inicial da aplicação.
            \item \textbf{Investments.jsx}: gerenciamento e visualização de investimentos.
            \item \textbf{Kpis.jsx}: exibição de indicadores e métricas financeiras.
            \item \textbf{Login.jsx}: autenticação e acesso dos usuários.
            \item \textbf{Releases.jsx}: gerenciamento e visualização de receitas e despesas financeiras.
        \end{itemize}
        
        \item \textbf{App.jsx}: componente principal da aplicação responsável pela renderização da interface.
        
        \item \textbf{App.css}: arquivo de estilos específicos do componente principal.
        
        \item \textbf{index.css}: arquivo global de estilização da aplicação.
        
        \item \textbf{main.jsx}: ponto de entrada da aplicação React, responsável por inicializar a renderização dos componentes.
    \end{itemize}

    \item \textbf{postcss.config.js}: configuração do PostCSS utilizado pelo Tailwind CSS.

\end{itemize}

\subsection{Organização Arquitetural}

A estrutura foi projetada para separar responsabilidades entre apresentação, gerenciamento de estado e configuração da aplicação. As páginas ficam concentradas no diretório \texttt{pages}, o compartilhamento de dados é realizado por meio do \texttt{FinanceContext}, enquanto os arquivos de configuração permanecem na raiz do projeto. Essa organização contribui para a legibilidade do código, facilita a manutenção e permite a evolução futura do sistema.


\section{Principais Classes e Funções}


Esta seção descreve a estrutura principal da aplicação de controle financeiro desenvolvida em React, detalhando os componentes, contextos e funções que compõem o sistema.

% ─────────────────────────────────────────────
\subsection{Ponto de Entrada — \texttt{main.jsx}}
% ─────────────────────────────────────────────

O arquivo \texttt{main.jsx} é o ponto de inicialização da aplicação. Ele é responsável por três responsabilidades centrais: configurar o roteamento, fornecer o contexto global de estado e montar a árvore de componentes no DOM.

\subsubsection*{Roteamento com \texttt{createBrowserRouter}}

O roteamento é gerenciado pela biblioteca \texttt{react-router-dom}. A função \texttt{createBrowserRouter} recebe um array de objetos de rota, cada um contendo os campos \texttt{path} (caminho da URL) e \texttt{element} (componente React a ser renderizado). As rotas definidas são:

\begin{itemize}
  \item \texttt{/} — renderiza o componente \texttt{Login};
  \item \texttt{home} — renderiza o componente \texttt{Home};
  \item \texttt{releases} — renderiza o componente \texttt{Releases};
  \item \texttt{kpis} — renderiza o componente \texttt{Kpis};
  \item \texttt{investments} — renderiza o componente \texttt{Investments}.
\end{itemize}

\subsubsection*{Montagem e Providers}

A montagem ocorre via \texttt{createRoot}, que recebe o elemento HTML com \texttt{id="root"} como âncora. O método \texttt{.render()} recebe a árvore de componentes envolvida por dois provedores:

\begin{itemize}
  \item \texttt{StrictMode} — ativa verificações adicionais de desenvolvimento fornecidas pelo próprio React, como detecção de efeitos colaterais inesperados;
  \item \texttt{FinanceProvider} — provedor do contexto global de finanças, que disponibiliza estado e funções para todos os componentes filhos;
  \item \texttt{RouterProvider} — injeta o roteador configurado em toda a aplicação.
\end{itemize}

% ─────────────────────────────────────────────
\subsection{Componente \texttt{Home}}
% ─────────────────────────────────────────────

O componente \texttt{Home} serve como painel principal da aplicação, exibindo um resumo financeiro do usuário e permitindo o acréscimo de saldo.

\subsubsection*{Consumo de Contexto}

Por meio do hook \texttt{useContext(FinanceContext)}, o componente acessa as seguintes variáveis e funções globais:

\begin{itemize}
  \item \texttt{saldoFinal} — valor numérico que representa o saldo disponível após descontar despesas e investimentos;
  \item \texttt{totalDespesas} — soma acumulada de todos os gastos registrados;
  \item \texttt{totalInvestido} — soma acumulada de todos os investimentos registrados;
  \item \texttt{adicionarSaldo} — função que incrementa o saldo global com um valor fornecido pelo usuário.
\end{itemize}

\subsubsection*{Estado Local}

O estado local \texttt{novoSaldo}, inicializado em string vazia via \texttt{useState}, armazena temporariamente o valor digitado no campo de entrada antes de ser submetido.

\subsubsection*{Função \texttt{handleAumentarSaldo}}

Esta função é acionada no envio do formulário. Ela realiza as seguintes etapas:

\begin{enumerate}
  \item Chama \texttt{e.preventDefault()} para impedir o recarregamento da página;
  \item Valida se \texttt{novoSaldo} não está vazio, é um número válido e é positivo — caso contrário, aborta a execução;
  \item Invoca \texttt{adicionarSaldo(novoSaldo)}, propagando o valor para o contexto global;
  \item Limpa o campo de entrada atribuindo string vazia ao estado \texttt{novoSaldo}.
\end{enumerate}

% ─────────────────────────────────────────────
\subsection{Componente \texttt{Releases}}
% ─────────────────────────────────────────────

O componente \texttt{Releases} é responsável pelo registro de despesas. O usuário informa uma descrição, um valor e uma categoria, e o lançamento é persistido no contexto global.

\subsubsection*{Consumo de Contexto}

Via \texttt{useContext(FinanceContext)}, o componente obtém:

\begin{itemize}
  \item \texttt{expenses} — array com todas as despesas registradas;
  \item \texttt{addExpense} — função que adiciona um novo objeto de despesa ao estado global.
\end{itemize}

\subsubsection*{Estados Locais}

Três estados locais controlam os campos do formulário:

\begin{itemize}
  \item \texttt{description} — texto descritivo do gasto;
  \item \texttt{value} — valor monetário em formato string (convertido para \texttt{float} no envio);
  \item \texttt{category} — categoria do gasto, com valor padrão \texttt{"Alimentação"}.
\end{itemize}

\subsubsection*{Função \texttt{handleSubmit}}

Ao enviar o formulário, a função executa:

\begin{enumerate}
  \item \texttt{e.preventDefault()} para suprimir o comportamento padrão do navegador;
  \item Validação mínima: se \texttt{description} ou \texttt{value} estiverem vazios, a execução é interrompida;
  \item Criação de um objeto de despesa com os campos \texttt{id} (gerado por \texttt{Date.now()}), \texttt{description}, \texttt{value} (convertido com \texttt{parseFloat}), \texttt{category} e \texttt{date} (formatada para o padrão brasileiro via \texttt{toLocaleDateString('pt-BR')});
  \item Chamada de \texttt{addExpense} com o objeto construído;
  \item Limpeza dos campos \texttt{description} e \texttt{value}.
\end{enumerate}

% ─────────────────────────────────────────────
\subsection{Componente \texttt{Kpis}}
% ─────────────────────────────────────────────

O componente \texttt{Kpis} (Key Performance Indicators) exibe métricas e indicadores visuais de gastos agrupados por categoria.

\subsubsection*{Consumo de Contexto}

Apenas \texttt{expenses} é consumido do contexto, pois o componente é exclusivamente de leitura e visualização.

\subsubsection*{Cálculo dos Totais por Categoria}

As categorias válidas são definidas estaticamente no array \texttt{categoriasValidas}:

\begin{center}
  \texttt{['Alimentação', 'Transporte', 'Lazer', 'Moradia', 'Outros']}
\end{center}

O objeto \texttt{totaisPorCategoria} é inicializado com \texttt{reduce}, atribuindo zero a cada categoria. Em seguida, o array \texttt{expenses} é percorrido com \texttt{forEach}: para cada despesa, o valor numérico é somado ao total geral (\texttt{totalCalculadoGeral}) e ao total da categoria correspondente. Despesas com categorias não reconhecidas são acumuladas em \texttt{'Outros'}.

\subsubsection*{Mapeamento de Cores}

O objeto \texttt{coresBarras} associa cada categoria a uma classe CSS do Tailwind, utilizada para colorir as barras de progresso exibidas na interface:

\begin{table}[h]
\centering
\begin{tabular}{ll}
\hline
\textbf{Categoria} & \textbf{Classe CSS} \\
\hline
Alimentação & \texttt{bg-amber-500} \\
Transporte  & \texttt{bg-blue-500} \\
Lazer       & \texttt{bg-purple-500} \\
Moradia     & \texttt{bg-pink-500} \\
Outros      & \texttt{bg-slate-500} \\
\hline
\end{tabular}
\caption{Mapeamento de categorias para classes de cor (Tailwind CSS)}
\end{table}

% ─────────────────────────────────────────────
\subsection{Componente \texttt{Investments}}
% ─────────────────────────────────────────────

O componente \texttt{Investments} segue a mesma estrutura do componente \texttt{Releases}, mas voltado ao registro de investimentos.

\subsubsection*{Consumo de Contexto}

Via \texttt{useContext(FinanceContext)}, o componente obtém:

\begin{itemize}
  \item \texttt{investments} — array com todos os investimentos registrados;
  \item \texttt{addInvestment} — função que insere um novo investimento no estado global.
\end{itemize}

\subsubsection*{Estados Locais}

Três estados controlam o formulário:

\begin{itemize}
  \item \texttt{description} — nome ou descrição do investimento;
  \item \texttt{value} — valor aportado;
  \item \texttt{type} — tipo do investimento, com valor padrão \texttt{"Renda Fixa"}.
\end{itemize}

\subsubsection*{Função \texttt{handleSubmit}}

O comportamento é análogo ao de \texttt{Releases}: após validação, um objeto com os campos \texttt{id}, \texttt{description}, \texttt{value}, \texttt{type} e \texttt{date} é construído e passado para \texttt{addInvestment}, seguido da limpeza dos campos do formulário.

% ─────────────────────────────────────────────
\subsection{Contexto Global — \texttt{FinanceContext}}
% ─────────────────────────────────────────────

O \texttt{FinanceContext} centraliza o estado financeiro da aplicação, seguindo o padrão \textit{Context API} do React. O provedor \texttt{FinanceProvider}, declarado em \texttt{main.jsx} como ancestral de toda a árvore de componentes, disponibiliza os seguintes dados e funções:

\begin{itemize}
  \item \textbf{Estado de despesas:} \texttt{expenses} e \texttt{addExpense};
  \item \textbf{Estado de investimentos:} \texttt{investments} e \texttt{addInvestment};
  \item \textbf{Saldo e derivados:} \texttt{saldoFinal}, \texttt{totalDespesas}, \texttt{totalInvestido} e \texttt{adicionarSaldo}.
\end{itemize}

Essa abordagem elimina a necessidade de \textit{prop drilling} — passagem explícita de dados por múltiplos níveis de componentes — tornando o estado acessível diretamente em qualquer folha da árvore via \texttt{useContext}.


\section{Fluxo de Execução da Aplicação}

Esta seção descreve, em ordem cronológica, como a aplicação é inicializada, como o estado global é construído, como o usuário navega entre as telas e como cada interação produz efeitos visíveis na interface. O objetivo é oferecer uma visão sistêmica do comportamento do código em tempo de execução.

% ─────────────────────────────────────────────
\subsection{Fase 1 — Inicialização}
% ─────────────────────────────────────────────

Quando o navegador carrega a aplicação, o ponto de entrada \texttt{main.jsx} é executado primeiro. Essa fase compreende três etapas sequenciais:

\begin{enumerate}
  \item \textbf{Construção do roteador.} A função \texttt{createBrowserRouter} processa o array de rotas e produz um objeto de roteador. Nenhum componente de página é renderizado ainda; o roteador apenas mapeia caminhos de URL a elementos React.

  \item \textbf{Criação da raiz React.} \texttt{createRoot} localiza o elemento \texttt{<div id="root">} no HTML e prepara o ponto de montagem. A partir deste momento, o React assume o controle do DOM dentro desse elemento.

  \item \textbf{Renderização inicial da árvore.} O método \texttt{.render()} é chamado com a seguinte hierarquia de componentes:

\begin{verbatim}
StrictMode
  └── FinanceProvider
        └── RouterProvider
\end{verbatim}

  O \texttt{FinanceProvider} é instanciado antes do \texttt{RouterProvider}, o que garante que o contexto financeiro já existe no momento em que qualquer componente de página for montado.
\end{enumerate}

% ─────────────────────────────────────────────
\subsection{Fase 2 — Construção do Estado Global}
% ─────────────────────────────────────────────

Com a montagem do \texttt{FinanceProvider}, o estado global da aplicação é inicializado. Internamente, o provedor define com \texttt{useState} os arrays \texttt{expenses} e \texttt{investments} (ambos começando vazios) e os valores numéricos de saldo. As funções \texttt{addExpense}, \texttt{addInvestment} e \texttt{adicionarSaldo} são criadas neste momento e permanecem estáveis durante todo o ciclo de vida da aplicação.

Todos esses valores são disponibilizados pelo \texttt{FinanceContext.Provider}, que os injeta na árvore de componentes. A partir deste ponto, qualquer componente filho pode consumir o contexto via \texttt{useContext(FinanceContext)} sem receber dados por \textit{props}.

% ─────────────────────────────────────────────
\subsection{Fase 3 — Resolução da Rota Inicial}
% ─────────────────────────────────────────────

Com o \texttt{RouterProvider} montado, o roteador lê a URL atual do navegador e determina qual componente renderizar. Na primeira visita, a URL corresponde ao caminho \texttt{"/"}, o que faz o roteador montar o componente \texttt{Login}. Nenhum outro componente de página existe no DOM neste instante.

A tela de \texttt{Login} é puramente de autenticação e não consome o \texttt{FinanceContext}. Ao completar o login, a navegação programática (tipicamente via \texttt{useNavigate} ou via componente \texttt{<Link>}) altera a URL para \texttt{/home}, disparando a fase seguinte.

% ─────────────────────────────────────────────
\subsection{Fase 4 — Renderização do Painel Principal (\texttt{Home})}
% ─────────────────────────────────────────────

Quando a URL muda para \texttt{/home}, o roteador desmonta o componente \texttt{Login} e monta o componente \texttt{Home}. Durante a montagem, ocorre a seguinte sequência interna:

\begin{enumerate}
  \item \textbf{Consumo do contexto.} O hook \texttt{useContext(FinanceContext)} é executado, e o componente recebe \texttt{saldoFinal}, \texttt{totalDespesas}, \texttt{totalInvestido} e \texttt{adicionarSaldo} diretamente do provedor.

  \item \textbf{Inicialização do estado local.} O hook \texttt{useState('')} cria o estado \texttt{novoSaldo}, que controla o campo de entrada de saldo.

  \item \textbf{Renderização da interface.} O componente produz o JSX com os valores atuais do contexto — neste momento todos zerados — e exibe o formulário de adição de saldo.
\end{enumerate}

\subsubsection*{Ciclo de Adição de Saldo}

Quando o usuário preenche o campo e submete o formulário, o seguinte fluxo ocorre:

\begin{enumerate}
  \item O evento \texttt{onSubmit} dispara \texttt{handleAumentarSaldo};
  \item A função valida o valor digitado em \texttt{novoSaldo};
  \item \texttt{adicionarSaldo(novoSaldo)} é chamada, atualizando o estado global no \texttt{FinanceProvider};
  \item O React detecta a mudança de estado no provedor e re-renderiza todos os componentes que consomem \texttt{FinanceContext}, incluindo o próprio \texttt{Home};
  \item A interface exibe o novo valor de \texttt{saldoFinal} imediatamente, sem recarregar a página;
  \item O estado local \texttt{novoSaldo} é limpo, esvaziando o campo de entrada.
\end{enumerate}

% ─────────────────────────────────────────────
\subsection{Fase 5 — Registro de Despesas (\texttt{Releases})}
% ─────────────────────────────────────────────

Ao navegar para \texttt{/releases}, o roteador monta o componente \texttt{Releases}. Três estados locais independentes — \texttt{description}, \texttt{value} e \texttt{category} — são inicializados. O estado global \texttt{expenses}, vindo do contexto, é utilizado para listar as despesas já registradas na tela.

\subsubsection*{Ciclo de Registro de Despesa}

\begin{enumerate}
  \item O usuário preenche os campos do formulário; cada digitação atualiza o respectivo estado local via \texttt{onChange};
  \item Ao submeter, \texttt{handleSubmit} é chamada;
  \item Um novo objeto de despesa é construído com \texttt{id} gerado por \texttt{Date.now()}, os valores dos estados locais e a data atual formatada;
  \item \texttt{addExpense(objeto)} é chamada, adicionando a despesa ao array \texttt{expenses} no \texttt{FinanceProvider};
  \item O React re-renderiza todos os consumidores do contexto: \texttt{Releases} exibe a nova despesa na lista, e \texttt{Home} (caso esteja montado) teria seu \texttt{totalDespesas} e \texttt{saldoFinal} atualizados;
  \item Os campos \texttt{description} e \texttt{value} são limpos; \texttt{category} mantém seu valor para facilitar lançamentos consecutivos da mesma categoria.
\end{enumerate}

% ─────────────────────────────────────────────
\subsection{Fase 6 — Registro de Investimentos (\texttt{Investments})}
% ─────────────────────────────────────────────

O fluxo do componente \texttt{Investments} é estruturalmente idêntico ao de \texttt{Releases}. A diferença reside nos dados manipulados: o campo \texttt{type} substitui \texttt{category}, o array \texttt{investments} substitui \texttt{expenses}, e a função \texttt{addInvestment} é invocada no lugar de \texttt{addExpense}.

O efeito cascata no contexto é análogo: após \texttt{addInvestment} ser chamada, o estado \texttt{totalInvestido} no provedor é recalculado, e o \texttt{saldoFinal} exibido em \texttt{Home} é atualizado proporcionalmente.

% ─────────────────────────────────────────────
\subsection{Fase 7 — Visualização de Indicadores (\texttt{Kpis})}
% ─────────────────────────────────────────────

O componente \texttt{Kpis} não possui estado local nem realiza mutações no contexto. É um componente exclusivamente de leitura. Ao ser montado via \texttt{/kpis}, o seguinte processamento ocorre:

\begin{enumerate}
  \item O array \texttt{expenses} é consumido do contexto em seu estado atual;
  \item O \texttt{reduce} inicializa \texttt{totaisPorCategoria} com zero para cada uma das cinco categorias válidas;
  \item O \texttt{forEach} itera sobre \texttt{expenses}, acumulando valores por categoria e no total geral;
  \item O JSX é gerado com os totais calculados, exibindo barras de progresso proporcionais ao percentual de cada categoria sobre o total;
  \item Como \texttt{Kpis} é consumidor do contexto, qualquer adição futura de despesa em \texttt{Releases} (em outra aba, por exemplo) causaria re-renderização automática e atualização dos indicadores.
\end{enumerate}

\section{Conclusão}

O presente trabalho apresentou o desenvolvimento de uma aplicação web de controle financeiro pessoal construída com React, tendo como eixo central o gerenciamento de estado por meio da Context API. Ao longo das seções anteriores, foram detalhados a arquitetura dos componentes, as responsabilidades de cada módulo e o fluxo de execução que governa o comportamento do sistema em tempo de execução.

A análise do projeto evidencia que a separação entre estado global e estado local foi a principal decisão arquitetural da aplicação. O \texttt{FinanceProvider} centraliza todos os dados financeiros relevantes — despesas, investimentos e saldo — e os disponibiliza de forma uniforme para os componentes consumidores sem a necessidade de passagem explícita de \textit{props}. Essa abordagem confere coesão ao sistema: qualquer alteração em um componente se propaga automaticamente para todos os demais que dependem do mesmo contexto, mantendo a interface consistente sem lógica de sincronização manual.

A estrutura de roteamento adotada com \texttt{react-router-dom} complementa essa organização ao isolar cada funcionalidade em uma rota dedicada. O componente \texttt{Home} fornece a visão consolidada do patrimônio; \texttt{Releases} e \texttt{Investments} encapsulam os formulários de entrada de dados; e \texttt{Kpis} oferece uma camada analítica sobre os lançamentos registrados. Essa divisão respeita o princípio de responsabilidade única e facilita tanto a manutenção quanto a extensão futura do sistema.

Do ponto de vista do fluxo de dados, a aplicação segue rigorosamente o modelo unidirecional característico do React: eventos de usuário disparam funções do contexto, que atualizam o estado global, que por sua vez provocam re-renderizações nos componentes afetados. Esse ciclo previsível reduz a superfície de erros e torna o comportamento do sistema rastreável em cada etapa.

Como perspectivas de evolução, o projeto poderia incorporar persistência de dados via \textit{localStorage} ou integração com uma API externa, de modo que as informações financeiras sobrevivam ao encerramento da sessão do navegador. A adição de autenticação real no componente \texttt{Login}, a implementação de filtros e intervalos de data nos lançamentos, e a exportação dos dados em formato PDF ou CSV são extensões naturais que seguiriam a mesma arquitetura estabelecida sem exigir reestruturação significativa do código existente.

Em síntese, a aplicação demonstra que os recursos nativos do React — Context API, hooks e roteamento baseado em componentes — são suficientes para construir um sistema de informação funcional, organizado e de fácil compreensão, sem a necessidade de bibliotecas de gerenciamento de estado externas. A clareza do fluxo de execução e a separação de responsabilidades entre os módulos constituem os principais pontos fortes da solução desenvolvida.

\clearpage
\begin{thebibliography}{99}

\bibitem{react}
META PLATFORMS, INC.
\textit{React Documentation}.
Disponível em: \url{https://react.dev/}.
Acesso em: 09 mai. 2026.

\bibitem{nodejs}
OPENJS FOUNDATION.
\textit{Node.js Documentation}.
Disponível em: \url{https://nodejs.org/docs/latest/api/}.
Acesso em: 09 mai. 2026.

\bibitem{tailwind}
TAILWIND LABS.
\textit{Tailwind CSS Documentation}.
Disponível em: \url{https://tailwindcss.com/docs}.
Acesso em: 11 mai. 2026.

\bibitem{postcss}
POSTCSS.
\textit{PostCSS Documentation}.
Disponível em: \url{https://postcss.org/}.
Acesso em: 12 mai. 2026.

\bibitem{npm}
NPM, INC.
\textit{npm Documentation}.
Disponível em: \url{https://docs.npmjs.com/}.
Acesso em: 09 mai. 2026.

\end{thebibliography}

\end{document}